\DocumentMetadata{
  pdfversion=2.0,
  pdfstandard=ua-2,
  lang=en-US,
  tagging=on,
  tagging-setup={math/setup=mathml-SE,tabsorder=structure}
}
\documentclass[11pt,letterpaper]{article}
\usepackage[margin=0.68in,includefoot]{geometry}
\usepackage{newcomputermodern}
\usepackage{amsmath,amscd,mathtools}
\providecommand{\square}{\mdlgwhtsquare}
\usepackage[hidelinks]{hyperref}
\hypersetup{
  pdftitle={Combinatorics PhD Exam, August 2018},
  pdfauthor={University of Florida Department of Mathematics},
  pdfsubject={Graduate mathematics examination},
  pdfdisplaydoctitle=true,
  bookmarksopen=true
}
\setcounter{secnumdepth}{0}
\setlength{\parindent}{0pt}
\setlength{\parskip}{0.28em}
\setlength{\emergencystretch}{3em}
\setlength{\leftmargini}{2.15em}
\setlength{\leftmarginii}{2.65em}
\setlength{\leftmarginiii}{3em}
\renewcommand{\labelenumi}{\textbf{\theenumi.}}
\pagestyle{empty}
\raggedbottom
\allowdisplaybreaks
\newenvironment{examproblems}[1]{%
  \begin{enumerate}%
  \setcounter{enumi}{#1}%
  \setlength{\topsep}{0.35em}%
  \setlength{\partopsep}{0pt}%
  \setlength{\itemsep}{0.48em}%
  \setlength{\parsep}{0.18em}%
}{\end{enumerate}}
\newenvironment{examsubparts}{%
  \begin{enumerate}%
  \setlength{\topsep}{0.18em}%
  \setlength{\partopsep}{0pt}%
  \setlength{\itemsep}{0.16em}%
  \setlength{\parsep}{0.1em}%
}{\end{enumerate}}
\newenvironment{examinstructions}{%
  \par\begingroup\itshape
}{%
  \par\endgroup\vspace{0.18em}
}
\makeatletter
\renewcommand\section{\@startsection{section}{1}{0pt}%
  {-1.25ex plus -.25ex minus -.1ex}%
  {0.55ex plus .1ex}%
  {\normalfont\large\bfseries}}
\renewcommand{\@maketitle}{%
  \newpage\null\vskip -1.2em%
  {\footnotesize\bfseries
    \href{https://www.ufl.edu/}{University of Florida}, \href{https://math.ufl.edu/}{Department of Mathematics}\par}%
  \vskip 0.18em%
  \tagstructbegin{tag=Title}%
    {\LARGE\bfseries\href{https://gma.math.ufl.edu/exams/combinatorics-phd/}{Combinatorics PhD Exam}, August 2018\par}%
  \tagstructend%
  \vskip 0.35em%
  \tagmcbegin{artifact}%
    \hrule height 1pt%
  \tagmcend%
  \vskip 0.55em%
}
\makeatother
\title{Combinatorics PhD Exam, August 2018}
\author{}
\date{}
\begin{document}
\maketitle
\thispagestyle{empty}
\begin{examproblems}{0}
\item
Find a generating function for the number of positive integer solutions to \(x_1+2x_2+3x_3+4x_4+5x_5=n\).
\item
Tickets numbered \(\{1,2,\ldots,n\}\) are drawn from an urn at random one after the other without replacement. What is the probability that the number~\(r\) is drawn on the \(k^{th}\) drawing?

\par
Among a population of \(n+1\) people, a rumor is spread at random. One person tells the rumor to a second, who in turn repeats it to a third person, etc. What is the probability that the rumor will be told~\(k\) times without being repeated to any person?
\item
Prove that there exists no simple, bipartite, planar graph with minimum degree at least~\(4\).
\item
Let~\(\mu\) be the largest eigenvalue of the adjacency matrix of a graph~\(G\) and~\(\Delta\) the maximum degree of~\(G\). Prove that \(\mu\leq\Delta\).

\par
(Hint: Consider the largest, in absolute value, coordinate of a corresponding eigenvector.)
\item
Show that a~\(k\)-regular graph of girth~\(4\) has at least \(2k\) vertices, and a~\(k\)-regular graph of girth~\(5\) has at least \(k^2+1\) vertices. Draw a~\(3\)-regular graph of girth~\(4\) with \(2\cdot 3=6\) vertices and a~\(3\)-regular graph of girth~\(5\) with \(3^2+1=10\) vertices.
\item
Let \(f_n\) be the number of all compositions of the integer~\(n\) into odd parts. Find the exponential order of the numbers \(f_n\).
\item
Recall that a permutation is called even if it has an even number of even cycles, and it is called odd if it has an odd number of even cycles. Also recall that a derangement is a permutation with no~\(1\)-cycles.

\par
Let \(A_n\) be the number of all even derangements of length~\(n\), and let \(B_n\) be the number of all odd derangements of length~\(n\). Find an exact formula for \(A_n-B_n\).
\item
This problem has two parts.
\begin{examsubparts}
\item[\textnormal{(a)}]
Prove that if a \((v,k,\lambda)\) BIBD exists with block set~\(\mathcal{B}\) on a point set~\(V\), then the set \(\mathcal{B}'=\{V\setminus B\mid B\in\mathcal{B}\}\) is a BIBD.
\item[\textnormal{(b)}]
Construct a \((7,4,2)\) BIBD.
\end{examsubparts}
\end{examproblems}
\end{document}
